\cotefile{32}
\scanpage{37}
Pour tout carré cartésien de schémas
\diagram{\begin{tikzcd}[ampersand replacement=\&]
 X' \arrow[r,"g'"] \arrow[d,"f'"] \& X \arrow[d,"f"]\\
 S' \arrow[r,"g"] \& S
\end{tikzcd}}
avec $f$ compactifiable, \deleted{\unsure{$\mathcal C$|$\ell$|$\mathcal E$}} et tout faisceau \inserted{d'anneaux} de torsion $\mathcal A$ sur $S$, l'isomorphisme de changement de base $\overline{\underline{\mathrm{XVII}}}$ 5.2
\margin{\unsure{$p$}}
\[
 g^* Rf_! \xrightarrow{\sim} Rf'_! g'^*
\]
entre foncteurs de $D(X,f^*\mathcal A)$ dans $D(S',g^*\mathcal A)$ se transpose en un isomorphisme
\[
 Rg'_* Rf'{}^! \xrightarrow{\sim} Rf^! Rg_* \tag*{(3.1.6.3)}
\]
entre foncteurs de $D^+(S',g^*\mathcal A)$ dans $D^+(X,f^*\mathcal A)$.

On obtient ainsi quatre ``\uline{flèches de changement de base}'' :
\diagram{\begin{tikzcd}[ampersand replacement=\&,column sep=large,row sep=large,
 execute at end picture={\draw[rounded corners,->] ([xshift=5mm,yshift=4mm]\tikzcdmatrixname-2-1.north east) -- ++(0,-3mm) -| ([xshift=-3mm]\tikzcdmatrixname-2-1.west) |- ([xshift=4mm,yshift=-3mm]\tikzcdmatrixname-3-1.south east) -- ++(0,-3mm);
 \draw[->] (\tikzcdmatrixname-1-1.south west) to[bend right=25] ([xshift=-6mm,yshift=-2mm]\tikzcdmatrixname-2-1.west);}]
 \tikz[baseline=(n.base)]{\node(n){$(3.1.6.4)$};\draw (n.center) ellipse (9mm and 3mm);} \& {} \& {} \& {}\\
 \deleted{\unsure{R}}g^*Rf_* \arrow[r] \& Rf'_*g'^* \& g^*Rf_! \arrow[r,"\sim", "\deleted{\leftarrow}"' near start] \& Rf'_!g'^*\\
 Rg^!Rf_* \& Rf'_*Rg'{}^! \arrow[l,"\sim"',"\deleted{\rightarrow}" near start] \& Rg^!Rf_! \& Rf'_!Rg'{}^! \arrow[l]
\end{tikzcd}}
Les flèches de la diagonale principale sont autoadjointes, et définies pour tout carré commutatif dès que leurs sources et buts ont un sens, \inserted{et en tous cas quand on considère les foncteurs envisagés comme définis sur \inserted{des catégories $D^+$.}} les isomorphismes de la seconde diagonale sont adjoints l'un de l'autre, et définis pour un carré \deleted{commutatif}\inserted{cartésien} de \deleted{\unsure{petits}} schémas, avec $f$ (resp $g$) compactifiable.

\inserted{Chacune de} ces flèches vérifie\deleted{\unsure{nt}} \deleted{des}\inserted{deux} compatibilités du type $\overline{\underline{\mathrm{XII}}}$ 4.4 pour $f$ ou $g$ composé de deux morphismes.
% slanted
\margin{Peut-on dire la raison?}
\margin{$g'^*Rf^!\xrightarrow{\sim}Rf'{}^!g^*$}
% slanted
\margin{\unsure{défini|dû} sans hyp. de finité}
% faint pencil at the foot, with a line between the two groups
\margin{(\unsure{est-ce que Verdier} \unsure{utilise} \unsure{pour calculer} $Rf^!$ \unsure{son} \illegible{1} $\longleftarrow$)}
\margin{(\unsure{manque} \illegible{1} : $\underline g$ loc. acyclique (p. ex. \uline{lisse}))}
% archive sheet number: 36
\scanpage{38}
\deleted{\uline{3.1.3 Construction du foncteur $Rf^!$}}
% archive sheet number: 37
\scanpage{39}
% slanted
\margin{Qui sont $\mathcal A,\mathcal B$ ?}
\deleted{\unsure{$\ell$}} 3.1.7. Pour $\underline F$ \deleted{complexe des}\inserted{un} $f^*\mathcal B$-Module\deleted{s} à droite sur $X$ et $\underline G$ \deleted{le complexe \unsure{de}} un $\mathcal A,\mathcal B$-Bimodule à gauche\inserted{sur $S$}, on dispose de
\[
 f_!^\bullet\underline F\otimes_{\mathcal B}\underline G
 \longrightarrow f_!^\bullet(\underline F\otimes_{\mathcal B} f^*\underline G).
 \tag*{(3.1.7.1)}
\]
Pour $\underline H$ faisceau de $\mathcal A$-modules sur $S$, cette flèche se transpose en
\[
 \underline{\mathrm{Hom}}_{\mathcal A}(f^*\underline G,f_\bullet^!\underline H)
 \longrightarrow f_\bullet^!\underline{\mathrm{Hom}}_{\mathcal A}(\underline G,\underline H)
 \tag*{(3.1.7.2)}
\]
Cette flèche se dérive en une flèche \deleted{d'induction}
\[
 R\underline{\mathrm{Hom}}_{\mathcal A}(f^*\underline{\deleted{\unsure{G}}L},Rf^!\underline{\deleted{\unsure{H}}M})
 \xrightarrow{\sim} Rf^!R\underline{\mathrm{Hom}}_{\mathcal A}(\underline{\deleted{\unsure{G}}L},\underline{\deleted{\unsure{H}}M}),
 \tag*{(3.1.7.3)}
\]
\margin{\tikz{\draw (.07,0) -- (0,.07) -- (.08,.43); \draw (0,.07) -- (-.02,0);}}
\inserted{($L\in\mathrm{ob}$ ?, $M\in\mathrm{ob}$ ?), appelée \uline{homomorphisme} \uline{d'induction}, dont on va vérifier que c'est un isomorphisme. Il} \deleted{qui} rend commutatif le diagramme suivant, pour $K\in D^-(X,f^*\mathcal B)$, $L\in D^-(S,\mathcal A,\mathcal B)$
\margin{\unsure{réf. not.|réf.} ?}
\deleted{$\mathrm{Hom}(Rf_!K,R\underline{\mathrm{Hom}}_{\mathcal A}(f^*L,Rf^!M))$} et $M\in D^+(S,\mathcal A)$.
\diagram{\resizebox{\linewidth}{!}{\begin{tikzcd}[ampersand replacement=\&,column sep=1.3em,row sep=2.5em]
 \mathrm{Hom}(K,R\underline{\mathrm{Hom}}_{\mathcal A}(f^*L,Rf^!M))
 \arrow[r,"\text{\textcircled{\scriptsize 1}}"] \arrow[d,leftrightarrow,"\sim" description]
 \& \mathrm{Hom}(K,Rf^!R\underline{\mathrm{Hom}}_{\mathcal A}(L,M))
 \arrow[r,leftrightarrow,"\sim","\mathrm{adj.}"']
 \& \mathrm{Hom}(Rf_!K,R\underline{\mathrm{Hom}}_{\mathcal A}(L,M))
 \arrow[d,leftrightarrow,"\sim" description]\\
 \mathrm{Hom}(K\overset{\mathbb L}{\otimes}_{\mathcal B}f^*L,Rf^!M)
 \arrow[r,leftrightarrow,"\sim","\mathrm{adj}"']
 \& \mathrm{Hom}(Rf_!(K\overset{\unsure{L|4}}{\otimes}f^*L),M)
 \arrow[r,"\text{\textcircled{\scriptsize 2}}"]
 \& \mathrm{Hom}(Rf_!K\overset{\mathbb L}{\otimes}L,M)
\end{tikzcd}}}
% slanted
\margin{terminologie}
Dans ce diagramme, \textcircled{\scriptsize 2} (déduit de la flèche de \deleted{changement de base}\inserted{projection}) est un isomorphisme ($\overline{\underline{\mathrm{XVII}}}$ 5.2). La flèche \textcircled{\scriptsize 1} est donc un isomorphisme quel que soit $K$. On vérifie facilement que ceci implique que (3.1.7.3) est un isomorphisme. \inserted{On a obtenu}

\margin{Préciser données $f$ comp. $X\to S$, $\mathcal A$ $\mathcal B$ \unsure{(?)}}
\uline{Proposition} 3.1.8. \uline{La flèche} 3.1.7.3 \uline{est un isomorphisme} : \uline{on a la ``\uwave{formule d'induction}''} :
\[
 R\underline{\mathrm{Hom}}_{\mathcal A}(f^*L,Rf^!M)
 \xrightarrow{\sim} Rf^!R\mathrm{Hom}_{\mathcal A}(L,M).
\]
\inserted{pour $L\in\mathrm{ob}\qquad,\ M\in\mathrm{ob}\qquad$.}
% archive sheet number: 38
\scanpage{40}
3.1.9 Soient $\mathcal A$ et $\mathcal B$ \deleted{d'un} deux faisceaux d'anneaux de torsion sur $S$ et soit $u:\deleted{\unsure{\mathcal A}}\mathcal B\longrightarrow\deleted{\unsure{\mathcal B}}\mathcal A$ un homomorphisme. Prenons dans 3.1.8, pour $L$, le complexe de $\mathcal A,\mathcal B$-bimodules \deleted{$\mathcal B$} reduit à $\deleted{\unsure{\mathcal B}}\mathcal A$ placé en degré 0. Le foncteur
\[
 \deleted{\underline F}\ \mathrm{Mod}(S,\mathcal A)\longrightarrow\mathrm{Mod}(S,\mathcal B):
 \underline F\longmapsto\underline{\mathrm{Hom}}_{\mathcal A}(L,\underline F),
\]
est le foncteur de restriction des scalaires $u_{.}$ de $\mathcal A$ à $\mathcal B$. La flèche d'adjonction, est ici un isomorphisme de foncteurs rendant commutatif le diagramme suivant.
\diagram{\begin{tikzcd}[ampersand replacement=\&,row sep=large]
 D^+\deleted{\mathrm{Mod}}(X,f^*\mathcal A) \arrow[r,"u_{.}"]
 \& D^+\deleted{\mathrm{Mod}}(\deleted{\unsure{S}}X,f^*\mathcal B)\\
 D^+\deleted{\mathrm{Mod}}(S,\mathcal A) \arrow[r,"u_{.}"] \arrow[u,"Rf^!"']
 \& D^+\deleted{\mathrm{Mod}}(S,\mathcal B)\,. \arrow[u,"Rf^!"']
\end{tikzcd}}
On voit que le faisceau $\mathcal A$ ne joue qu'un rôle bidon dans la définition de $Rf^!$.
% slanted
\margin{\illegible{3} $\searrow$}
% archive sheet number: 39
\scanpage{41}
\deleted{Prenons} \deleted{Soit $\unsure{\mathcal A|\mathcal B}$}
% archive sheet number: 40
\scanpage{42}
% archive sheet number: 41a (upper slip)
% slanted
\margin{Le rôle de 3.1.13 est bien obscur. Ne peut-on le \unsure{vider|virer} avantageusement ?}
On a dans ce cas, posant $K^!=K'{}^!$,
\[
 K^{* !}\underline F=a\bigl(W\longmapsto\mathrm{Hom}_{V_W}(K(W),\underline F\mid V_W)\bigr),
\]
\[
 \underline H^n(RK^!L)_x\simeq\underset{\unsure{W(x)|W(\bar x)}}{\underrightarrow{\lim}}\,
 \mathrm{Ext}^n_{V_W}(K(W),L\mid V_W).
\]

% Lower sheet; a long pale oblique line crosses the writing from top to bottom.
\tikz[remember picture,overlay]{\coordinate (scanXLIIstart) at (1.3,0);}
% slanted
\margin{À quoi sert cette section 3.1.14, qu'y fait-on ?}
3.1.\deleted{\unsure{16|10}}14. Appliquons les considérations qui précèdent à un morphisme compactifiable de schémas $f:X\longrightarrow S$.
D'après les démonstrations de 3.1.\deleted{\unsure{11}}3 et 3.1.4, si $f=\bar f j$ est une compactification de $f$ et si les fibres de $f$ sont de dimension $\leq d$, le foncteur $Rf^!$ se construit par le procédé de 3.1.10, pour
\margin{3.1.10}
\[
 K(U)=\bar f_*\tau^{\unsure{\bullet|*}}_{\leq d}C^{*}_{\ell}\mathcal A_U=f_!^\bullet(\mathcal A_U).
\]
\margin{$\tau''_{\leq 2d}$ ?}
\inserted{où $A_U$ ---}
% slanted
\margin{dire qui est $A_U$}

Le corollaire 3.1.12 va permettre de donner d'autres procédés de calcul\unsure{/} du foncteur $Rf^!$. Supposons que $n\mathcal A=0$. On a alors ($\overline{\underline{\mathrm{XVII}}}$ 5.2\tikz[baseline]{\draw (0,-.07) -- (.06,.42); \draw (.02,.06) -- (.10,-.02);})
\margin{\tikz{\draw (0,-.07) -- (.06,.42); \draw (.02,.06) -- (.10,-.02);}}
\[
 Rf_!\mathcal A_U\simeq\mathcal A\overset{\mathbb L}{\otimes}Rf_!(\mathbb Z/n)_U.
\]
En fait, $f_!^\bullet((\mathbb Z/n)_U)$ est un complexe de faisceaux \uline{plats} : il suffit de le voir pour $n$\deleted{$=$}\inserted{de la forme} $\ell^k$ \deleted{primaire}, \deleted{\unsure{et}} la platitude \deleted{\unsure{degré par degré}} d'un faisceau $\underline F$ \deleted{\unsure{de gr}} de $\mathbb Z/\ell^k$-modules signifie que pour $i<k$, $\underline F/\ell\underline F\xrightarrow{\sim}\ell^i\underline F/\ell^{i+1}\underline F$, et ceci résulte pour les faisceaux $f_!^i(\mathbb Z/n)_U$ de l'exactitude de $f_!^i$ et de la platitude \deleted{\unsure{de $f_!^i$}} $(\mathbb Z/n)_U$.
\deleted{On a donc un quasi-isomorphisme} La flèche \inserted{canonique}
\[
 \deleted{f_!^\bullet(\mathbb Z/n)_U\xrightarrow{\sim}}
\]
\[
 f_!^\bullet(\mathbb Z/n)_U\otimes\mathcal A\xrightarrow{\sim}f_!^\bullet(\mathcal A_U)
\]
est donc un quasi-isomorphisme et on peut aussi bien décrire (3.1.12)
\tikz[remember picture,overlay]{\draw[gray,thin] (scanXLIIstart) -- (0,0);}
% archive sheet number: 41
\scanpage{43}
% A pale sweeping diagonal crosses the body of the sheet.
\tikz[remember picture,overlay]{\coordinate (scanXLIIIstart) at (1.2,0);}
le foncteur $Rf^!$ comme fourni par la \deleted{\unsure{correspondance|complexe}} donnée 3.1.10
\[
 U\longmapsto f_!^\bullet((\mathbb Z/n)_U)
\]
% slanted
\margin{$\otimes_{\mathbb Z/n}\mathcal A$ ? sinon, expliciter \unsure{de nouveau|à nouveau|le nouveau} $\mathcal A,\mathcal B$}
Rappelons\tikz[baseline]{\draw (0,-.06) -- (.18,.34); \draw (.04,.03) -- (.12,-.07);}\margin{réf ?} que la \uline{résolution plate canonique} \inserted{$L_\bullet(\underline F)$} d'un Module $\underline F$ s'obtient en représentant $\underline F$ comme quotient du Module engendré par le faisceau d'ensembles $\underline F$ : $\varepsilon:L_0(\underline F)\longrightarrow F$, en appliquant ensuite cette construction à $\mathrm{Ker}(\varepsilon)$, et en itérant le procédé.

Le foncteur $Rf^!$ est encore fourni, \inserted{par la méthode de 3.1.10,} par la donnée
\[
 U\longmapsto\inserted{K_0(U)\overset{\mathrm{dfn}}{=}}L_\bullet(f_!^\bullet(\mathbb Z/n)_U)\deleted{=K_0(U)}.
\]
Si maintenant $\underline G$ est un complexe de $\mathcal A$-Modules acycliques\unsure{,|.} $(\qquad\qquad)$, on a
\margin{??}
\[
 K_0^!(\underline G)\xrightarrow{\sim}RK_0^!(\underline G).
\]
On \deleted{en déduit}\inserted{retrouve} que le faisceau d'anneaux $\mathcal A$ ne joue qu'un rôle bidon dans la construction de $\deleted{\unsure{RK^!|RK_0^!}} Rf^!$, \deleted{on a}
\tikz[remember picture,overlay]{\draw[gray,thin] (scanXLIIIstart) to[bend left=4] (0,0);}
% archive sheet number: 42
\scanpage{44}
Je n'ai pas vérifié que \tikz[baseline,overlay]{\draw[gray] (-.1,.48) -- (.16,-.42);}cette construction \unsure{coïncide|coincide} avec celle de 3.1.9.
% archive sheet number: 43
\scanpage{45}
3.1.\deleted{\unsure{13}}15. On peut \deleted{\unsure{encore}} décrire $Rf^!$ par le procédé de 3.1.13. \deleted{\unsure{Résultat}} Soit en effet un morphisme compactifiable $f$ :
\diagram{\begin{tikzcd}[ampersand replacement=\&]
 X \arrow[r,hook,"j"'] \arrow[d,"f"'] \& \bar X \arrow[dl,"\bar f"]\\
 S \& {}
\end{tikzcd}}
et $\mathcal A$ un faisceau d'anneaux sur $S$ tel que $n\mathcal A=0$ $(n\geq1)$. Pour $W=(U,V,\unsure{\varphi|\psi})\in\mathrm{Ob}\,X_f$ (notation de 3.1.13), soit
\[
 K(W)=\bar f_{V*}\tau_{\leq 2d}C^*_{\ell,\bar X_V}(\mathcal A_U)
 \tag*{(3.1.15.1)}
\]
\margin{$\tau''$ ?}
% slanted
\margin{signification pas claire}
(résolution calculée sur $\bar X_V=\bar X\times_S V$). \deleted{On a un quasi-isomorphisme}
\[
 \deleted{j_!K(W)}
\]
% diagram in the left margin, with the diagonal arrows directed down and left
\diagram{\begin{tikzcd}[ampersand replacement=\&,row sep=small,column sep=small]
 \bar X \arrow[d] \arrow[dd,bend right=48,"\bar f"'] \& \bar X_V \arrow[d] \arrow[dd,bend right=42,"\bar f_V"'] \& {}\\
 X \arrow[d,"f"'] \& X_V \arrow[l] \arrow[d,"f_V"] \& U \arrow[l]\\
 S \& V \arrow[l,"j"'] \& {}
\end{tikzcd}}
Pour $j$ \deleted{inclusion de}\inserted{morphisme structurel} $V$ dans $S$, on a un quasi-isomorphisme
\[
 j_!K(W)\xrightarrow{\sim}\bar f_*\tau_{\leq 2d}C^*_{\ell,\bar X}(\mathcal A_U)
\]
\margin{$(j\bar f_V)_*$ ?}
($\overline{\underline{\mathrm{XVII}}}$ 5.2.\tikz[baseline]{\draw (0,-.07) -- (.10,.40); \draw (.02,.02) -- (.10,-.06);})
\margin{réf}
D'après 3.1.9,\margin{?} le système (3.1.13)\margin{?} des $K(W)$ permet donc de calculer $Rf^!$.
% slanted
\margin{pas clair}
\deleted{Ceci donne un sens à la ``\uline{formule d'induction}''}
% archive sheet number: 44
\scanpage{46}
\deleted{Proposition 3.1.12. Le morphis}
% archive sheet number: 45
\scanpage{47}
% slanted
\margin{\unsure{remonter|remonte|remarque} en cor. à 3.1.4 ?}
3.1.16. L'isomorphisme \deleted{d'induction}\inserted{d'adjonction} s'écrit
\[
 \mathrm{Hom}(Rf_!K,L)\simeq\mathrm{Hom}(K,Rf^!L)
 \tag*{(3.1.16.1)}
\]
\inserted{avec} \deleted{\{} $K\in\mathrm{Ob}\,D(X,f^*\mathcal A)$ et $L\in\mathrm{Ob}\,D^+(S,\mathcal A)$. \deleted{\}} Pour $K\in\mathrm{Ob}\,D^-(X,f^*\mathcal A)$ et $L\in\mathrm{Ob}\,D^+(S,f^*\mathcal A)$\ensuremath{{}_{*}} \deleted{cet isomorphisme} complexe d'injectifs, cet isomorphisme s'écrit, au niveau des complexes,
\margin{$L$, ou $K$ et $L$ ?}
\[
 \mathrm{Hom}^\bullet(f_!^\bullet K,L)\simeq
 \mathrm{Hom}^\bullet(K,f_\bullet^!L).
\]
Pour tout \deleted{$V\in\mathrm{Ob}\,S_{\mathrm{et}}$}\inserted{$j:V\to S$ dans $S_{\mathrm{et}}$}, on dispose de
\[
\begin{aligned}
 \mathrm{Hom}_{f^*V}^\bullet(K,f_\bullet^!L)
 &=\mathrm{Hom}^\bullet(j'_!j'^*K,f_\bullet^!L)
 \underset{\sim}{\longleftrightarrow}\mathrm{Hom}^\bullet(f_!^\bullet(j'_!j'^*K),L)\longrightarrow\\
 &\deleted{\longrightarrow\mathrm{Hom}\ \mathrm{Hom}}\\
 &\longrightarrow\mathrm{Hom}(j_!j^*(f_!^\bullet K),L)
 =\mathrm{Hom}_{\deleted{\unsure{S}}V}(f_!^\bullet K,L)
\end{aligned}
\]
d'où un morphisme
\[
 Rf_*R\underline{\mathrm{Hom}}(K,Rf^!L)\xrightarrow{\sim}
 R\underline{\mathrm{Hom}}(Rf_!K,L)
 \tag*{(3.1.16.2)}
\]
% margin note written sideways
\margin{devrait remonter en cor. 3.1.4.8 à 3.1.4 (conséquence triviale de 3.1.4.7)}
que par localisation de 3.1.16.1 on vérifie être un isomorphisme.

3.1.17. \uline{Exemple} :
% slanted
\margin{\unsure{Videz|Vider}, ou en faire \unsure{remarque|rem.|venir} d'une ligne.}
Soi\inserted{en}t $X$ un schéma compactifiable sur le spectre d'un corps algébriquement clos $k$ \deleted{\unsure{et}}, $x$ un point fermé de $X$ et $n$ un entier $\geq1$ inversible dans $k$\unsure{,|.} La formule (3.1.11.1) fournit, puisque $\mathbb Z/n\mathbb Z$ est un $\mathbb Z/n\mathbb Z$-module injectif ; \deleted{Pour $U(x)$ parcourant les voisinages étales de $x$}
\[
 \underline H^{-i}(Rf^!\mathbb Z/n\mathbb Z)_x
 \simeq\underset{U(x)}{\underrightarrow{\lim}}\,
 \mathrm{Hom}(\deleted{\unsure{\mathbb Z/}}H_c^i(U,\mathbb Z/n\mathbb Z),\mathbb Z/n\mathbb Z),
\]
\inserted{où $U(x)$ est la catégorie des voisinages étales de $x$ dans $X$.}
% slanted
\margin{pourquoi pas $W(x),W$ ?}
Il résulte vraisemblablement du théorème de bidualité (SGA 5 I \qquad) (qui dépend de la résolution des singularités)
\margin{?}
% archive sheet number: 46
\scanpage{48}
% archive sheet number: 47 (upper left)
\tikz[remember picture,overlay]{\coordinate (scanXLVIIIstart) at (-.25,.25);}que la flèche\unsure{/}
\[
 H_x^i(X,\mathbb Z/n\mathbb Z)\longrightarrow
 \deleted{\unsure{H^i}}\underset{U(x)}{\underleftarrow{\text{``}\lim\text{''}}}\,
 H_c^i(U,\mathbb Z/n\mathbb Z)
\]
est un isomorphisme de pro-objets, le second membre étant représenté par le premier. On aurait alors un isomorphisme
\[
 \underline H^{-i}(Rf^!\mathbb Z/n\mathbb Z)_x
 \xrightarrow{\sim} H_x^i(X,\mathbb Z/n\mathbb Z)^\vee
\]

\tikz[remember picture,overlay]{\draw[gray,thin] (scanXLVIIIstart) to[bend right=8] (-.25,.35);}
La propriété de ``finitude'' suivante résulte aussitôt \deleted{Pro} de \deleted{ce que}\inserted{la construction} des \unsure{foncteur|foncteurs} $f_i^!$ (3.1.4.5), \inserted{et du fait qu'ils} transforment injectifs en injectifs.

\uline{Proposition} 3.1.18. (i) \uline{\deleted{\unsure{Soit|Si}}Supp\inserted{osons} le morphisme compactifiable} $f:X\longrightarrow S$ \uline{est de dimension relative} $\leq d$, \deleted{\unsure{si alors}} \deleted{\unsure{alors}}\inserted{et soit} $L\in D^+(S,\mathcal A)$, \deleted{\uline{alors}}\inserted{tel que} \deleted{si} $\underline H^i(L)=0$ \uline{pour} $i\geq k$, \uline{alors} $\underline H^i(L)=0$ \uline{pour} $i\geq k-2d$
% slanted
\margin{devrait remonter avant 3.1.5 ! ou imméd. après}

\deleted{(ii) \uline{Si} $\dim\mathrm{\,inj}(L)\leq k$, \uline{alors dim} $\mathrm{inj}(Rf^!L)\leq \unsure{kd|k1}$.}

(ii) $\dim\mathrm{\,inj}(Rf^!L)\leq\dim\mathrm{\,inj}(L)$.

% slanted
\margin{remonter également avant 3.1.6. \uline{bref}}
3.1.19. Lorsque $f$ est quasi-fini et si on fait $d=0$ dans la construction 3.1.4 du foncteur $Rf^!$, on trouve \inserted{par 3.1.3} que le foncteur $f_!$ admet un adjoint à droite $f^!$, dont $Rf^!$ est le dérivé.
% slanted
\margin{de quoi dans quoi ?}

En particulier \uline{si} $f$ \uline{est une immersion} $f:X\hookrightarrow S$, on trouve que \uline{$Rf^!$ est le dérivé du foncteur} ``\uline{sections à support dans $X$}''. En particulier
\[
 R^if^!_{\deleted{\unsure{?}}}L=\underline H_X^i(L)
 \tag*{(3.1.19.1)}
\]
\margin{\unsure{Dire} aussi que si $j$ est \unsure{éta.|étale}, $Rj^!=j^*$ ; corollaire : ``nature locale'' de $Rf^!$}

Le plan de 3.1 semble flou après 3.1.4. On pourrait prendre

a) Propriétés formelles de $Rf^!$ pour $f$ fixé : \deleted{\unsure{a)}} 3.1.5, 3.1.18, 3.1.19, 3.1.8, 3.1.9\margin{?}

b) Propriétés formelles pour $f$ variable : 3.1.6 $+$ \fbox{\shortstack[l]{chgt de base par morphisme\\ lisse (ou loc. acyclique)}}

c) Sorites sur construction locale de $Rf^!$ : 3.1.10 \unsure{à|''} 3.1.13

\diagram{\begin{tikzcd}[ampersand replacement=\&,column sep=small,row sep=small]
 X \arrow[d,dash,"f"'] \& X\times_S X \arrow[l,"p_1"'] \arrow[d,"p_2"]
 \& X \arrow[l,"\Delta"']\\
 S \& X \arrow[l,dash,"f"] \& {}
\end{tikzcd}}
\[
 \begin{aligned}
 \deleted{f}f^!&\simeq(p_1\Delta)^*f^!\simeq\Delta^*(p_1^*f^!)
 \underset{\text{chgt de base par lisse}}{\simeq}\Delta^*(p_2^!f^*)=\deleted{\unsure{\Delta^!p_2^!f^*}}\\
 &\underset{\text{pureté relative ??}}{\simeq}\Delta^!p_2^!f^*[2d](d)\simeq f^{\unsure{!|*}}[2d](d)
 \end{aligned}
\]
\margin{$\uparrow$ \uline{à ne pas} \uline{oublier} !}
% The arrow from the last note points to the boxed base-change item in b).
